
%We perform a robustness studies to validate our results. 
We investigate if the results of our study hold up when we focus the analysis upon the trading hours of NYSE and compute one model per trading hour. By this, we investigate if seasonality of trading is impacting our results. As there are 7 trading hours, this leads to 7 models for each of the time series to be analysed, namely return, integrated volatility and jump volatility. %and trading volume. 
The results of these analysis are displayed in Tables \ref{table_hourly_RET}, \ref{table_hourly_IV} and \ref{table_hourly_JV}. %, and \ref{table_hourly_VOl}.
The results indicate that the change in posting behaviour of the VNs and ANVNs plays an even stronger role when focusing upon the individual trading hours. Across the hours and data series, the $R^2$ and adj. $R^2$ are much higher. For the return series, Table \ref{table_hourly_RET}, we observe that for all hours but the first and last one, a higher proportion of the variability is explained by the  change points of VN and ANVN compared to the network statistics. In all cases, considering the change points and network statistics improves the $R^2$ and adj. $R^2$ the most. Notably, the models strongly outperform a standard autoregressive model. For hour 5, the AR model only yields 0.1 adj. $R^2$ but the model considering AC and the change points of VN and ANVN yields 0.48 adj. $R^2$. 

In terms of modelling the IV and JV by the trading hours, similar results as for the entire time series are observed but with much higher $R^2$ and adj. $R^2$, Table \ref{table_hourly_IV}. For various hours, the change points of the VN and ANVN add so much explainability of the variability to the base model, that the network statistics only explain very little, consider hours 4, 5, 6, or 7. For hour 4 the change points of the VN and ANVN explain such a large proportion of the variability, $R^2 = 0.9$, whereas the model with the network statistics only explains $R^2 = 0.67$, that it can be concluded the change in behaviour of the users on WSB explain almost the entire variability of the IV. 
For the JV, Table \ref{table_hourly_JV}, we observed that the change points of the VN and ANVN add a lot more explainability to the variability of the time series in the per trading hour analysis compared to the entire time series. For the latter, modelling the JV yielded $R^2 = 0.129$ for the model consisting of change points, network statistics and activity. For the per trading hour analysis, they explain at least 41\% of the variability of the system, for hour 4 even 86\% of the variability are explained by the change points, network statistics and activity. The model containing only network activity and change points of VN and ANVN explain in most cases, but hour 1, 6, 7, more of the JVs variability than the network statistics. 
%For the trading volume, Table  \ref{table_hourly_VOl}, we observe similarly to the modelling of the entire time series, that the network statistics and changepoints do not explain the variability of it well. For hours 1-4, the adj. $R^2$ is even frequently negative. For the hours 5-7 the results are better and we observe that the adj. $R^2$ is better for the model containing only network activity and the changepoints of the VN and ANVN compared to the model with the network statistics. 

To sum up, these results show that the change in behaviour from the users on WSB, measured by the change points of the VN and ANVN, explain a large proportion of the variability of the return series, integrated variance and jump volatility. %and to some extend also the trading volume series. 
The results also show that during market trading hours, an even larger proportion of the affected stocks variability can be explained. The latter result is in line with our observation that the hour of the day and trading hours have an effect upon the WSB users posting behaviour, compare section 5. 

% latex table generated in R 4.2.1 by xtable 1.8-4 package
% Tue Feb  6 18:16:33 2024
\begin{table}[ht]
\centering
\caption{Return: This table shows that the $R^2$ and adj. $R^2$ for the return prediction models regressed onto the observations 
                     as by trading hours. Column (1) shows the results for the autoregressive benchmark model, (2)-(6) add variables to the benchmark model.  The results show that regardless of the trading hour, the change points of the VN and ANVN strongly 
                     improve the modelling of the return series. Notably the $R^2$ and adj. $R^2$ are much higher in the hourly regressions than 
                     in the models for the entire time series.} 
\label{table_hourly_RET}
\begingroup\scriptsize
\begin{tabular}{ll|rrrrrr}
   &   & (1) & (2) & (3) & (4) & (5) & (6) \\ 
  \hline
\multirow{2}{*}{Hour 1} & $R^2$ & 0.07 & 0.09 & 0.1 & 0.31 & 0.33 & 0.49 \\ 
    & adj. $R^2$ & 0.03 & -0.01 & -0.01 & 0.17 & 0.1 & 0.25 \\ 
   \hline
\multirow{2}{*}{Hour 2} & $R^2$ & 0.1 & 0.13 & 0.42 & 0.58 & 0.41 & 0.67 \\ 
    & adj. $R^2$ & 0.08 & 0.08 & 0.36 & 0.52 & 0.29 & 0.57 \\ 
   \hline
\multirow{2}{*}{Hour 3} & $R^2$ & 0.29 & 0.3 & 0.38 & 0.46 & 0.42 & 0.55 \\ 
    & adj. $R^2$ & 0.26 & 0.25 & 0.31 & 0.38 & 0.31 & 0.41 \\ 
   \hline
\multirow{2}{*}{Hour 4} & $R^2$ & 0.52 & 0.52 & 0.62 & 0.65 & 0.59 & 0.71 \\ 
    & adj. $R^2$ & 0.5 & 0.49 & 0.58 & 0.6 & 0.51 & 0.62 \\ 
   \hline
\multirow{2}{*}{Hour 5} & $R^2$ & 0.13 & 0.16 & 0.29 & 0.55 & 0.23 & 0.61 \\ 
    & adj. $R^2$ & 0.1 & 0.1 & 0.21 & 0.48 & 0.07 & 0.49 \\ 
   \hline
\multirow{2}{*}{Hour 6} & $R^2$ & 0.16 & 0.17 & 0.32 & 0.36 & 0.32 & 0.46 \\ 
    & adj. $R^2$ & 0.13 & 0.11 & 0.24 & 0.26 & 0.17 & 0.29 \\ 
   \hline
\multirow{2}{*}{Hour 7} & $R^2$ & 0.18 & 0.2 & 0.26 & 0.29 & 0.33 & 0.44 \\ 
    & adj. $R^2$ & 0.15 & 0.14 & 0.18 & 0.19 & 0.19 & 0.26 \\ 
   \hline & Vital nodes & No & No & Yes & Yes & No & Yes \\ & Non-vital nodes & No & No & No & Yes & No & Yes \\ & Network activity & No & Yes & Yes & Yes & Yes & Yes \\ & Network statistics & No & No & No & No & Yes & Yes \\ \hline
\hline
\end{tabular}
\endgroup
\end{table}

% latex table generated in R 4.2.1 by xtable 1.8-4 package
% Tue Feb  6 18:17:07 2024
\begin{table}[ht]
\centering
\caption{Integrated variance: This table shows the $R^2$ and adj. $R^2$ for the integrated variance prediction models regressed onto the observations 
                     as by trading hours. Column (1) shows the results for the autoregressive benchmark model, (2)-(6) add variables to the benchmark model. The results show that regardless of the trading hour, the change points of the VN and ANVN strongly 
                     improve the modelling of the integrated variance. Notably the $R^2$ and adj. $R^2$ are much higher in the hourly regressions than 
                     in the models for the entire time series. For Hour 4 we observe a particularly strong improvement from the VN and ANVN change points.} 
\label{table_hourly_IV}
\begingroup\scriptsize
\begin{tabular}{ll|rrrrrr}
   &   & (1) & (2) & (3) & (4) & (5) & (6) \\ 
  \hline
\multirow{2}{*}{Hour 1} & $R^2$ & 0.53 & 0.56 & 0.59 & 0.68 & 0.65 & 0.74 \\ 
    & adj. $R^2$ & 0.51 & 0.51 & 0.54 & 0.61 & 0.53 & 0.61 \\ 
   \hline
\multirow{2}{*}{Hour 2} & $R^2$ & 0.65 & 0.67 & 0.68 & 0.84 & 0.74 & 0.9 \\ 
    & adj. $R^2$ & 0.64 & 0.65 & 0.65 & 0.82 & 0.68 & 0.87 \\ 
   \hline
\multirow{2}{*}{Hour 3} & $R^2$ & 0.62 & 0.63 & 0.69 & 0.76 & 0.72 & 0.82 \\ 
    & adj. $R^2$ & 0.61 & 0.61 & 0.66 & 0.72 & 0.66 & 0.77 \\ 
   \hline
\multirow{2}{*}{Hour 4} & $R^2$ & 0.52 & 0.53 & 0.69 & 0.9 & 0.67 & 0.91 \\ 
    & adj. $R^2$ & 0.51 & 0.5 & 0.65 & 0.88 & 0.6 & 0.89 \\ 
   \hline
\multirow{2}{*}{Hour 5} & $R^2$ & 0.6 & 0.61 & 0.64 & 0.8 & 0.73 & 0.85 \\ 
    & adj. $R^2$ & 0.59 & 0.58 & 0.6 & 0.78 & 0.68 & 0.8 \\ 
   \hline
\multirow{2}{*}{Hour 6} & $R^2$ & 0.78 & 0.8 & 0.81 & 0.89 & 0.85 & 0.91 \\ 
    & adj. $R^2$ & 0.77 & 0.79 & 0.79 & 0.87 & 0.82 & 0.88 \\ 
   \hline
\multirow{2}{*}{Hour 7} & $R^2$ & 0.69 & 0.71 & 0.72 & 0.82 & 0.75 & 0.85 \\ 
    & adj. $R^2$ & 0.68 & 0.68 & 0.69 & 0.79 & 0.7 & 0.8 \\ 
   \hline & Vital nodes & No & No & Yes & Yes & No & Yes \\ & Non-vital nodes & No & No & No & Yes & No & Yes \\ & Network activity & No & Yes & Yes & Yes & Yes & Yes \\ & Network statistics & No & No & No & No & Yes & Yes \\ \hline
\hline
\end{tabular}
\endgroup
\end{table}

% latex table generated in R 4.2.1 by xtable 1.8-4 package
% Tue Feb  6 18:17:07 2024
\begin{table}[ht]
\centering
\caption{Jump volatility: This table shows the $R^2$ and adj. $R^2$ for the jump volatility prediction models regressed onto the observations 
                     as by trading hours. Column (1) shows the results for the benchmark model with network activity variables, (2)-(5) add variables to the benchmark model.  The results show that regardless of the trading hour, the change points of the VN and ANVN strongly 
                     improve the modelling of the jump volatility. We observe that contrary to the regression onto the entire time series, 
                     the change points of VN and ANVN strongly improve the modelling of the jump volatility and explains parts the network statistic cannot. 
                     This observation is particularly pronounced for Hours 2-5, so the morning and midday trading hours. Further, 
                     the $R^2$ and adj. $R^2$ are much higher in the hourly regressions than in the models for the entire time series.} 
\label{table_hourly_JV}
\begingroup\scriptsize
\begin{tabular}{ll|rrrrr}
   &   & (1) & (2) & (3) & (4) & (5) \\ 
  \hline
\multirow{2}{*}{Hour 1} & $R^2$ & 0.14 & 0.15 & 0.19 & 0.34 & 0.41 \\ 
    & adj. $R^2$ & 0.1 & 0.09 & 0.09 & 0.17 & 0.18 \\ 
   \hline
\multirow{2}{*}{Hour 2} & $R^2$ & 0.14 & 0.14 & 0.38 & 0.37 & 0.56 \\ 
    & adj. $R^2$ & 0.11 & 0.09 & 0.31 & 0.27 & 0.45 \\ 
   \hline
\multirow{2}{*}{Hour 3} & $R^2$ & 0.06 & 0.1 & 0.31 & 0.23 & 0.41 \\ 
    & adj. $R^2$ & 0.03 & 0.04 & 0.23 & 0.11 & 0.26 \\ 
   \hline
\multirow{2}{*}{Hour 4} & $R^2$ & 0.04 & 0.49 & 0.84 & 0.44 & 0.86 \\ 
    & adj. $R^2$ & 0 & 0.45 & 0.82 & 0.36 & 0.83 \\ 
   \hline
\multirow{2}{*}{Hour 5} & $R^2$ & 0.02 & 0.05 & 0.42 & 0.25 & 0.57 \\ 
    & adj. $R^2$ & -0.02 & -0.01 & 0.36 & 0.13 & 0.46 \\ 
   \hline
\multirow{2}{*}{Hour 6} & $R^2$ & 0.02 & 0.05 & 0.11 & 0.38 & 0.43 \\ 
    & adj. $R^2$ & -0.02 & -0.02 & 0.01 & 0.28 & 0.28 \\ 
   \hline
\multirow{2}{*}{Hour 7} & $R^2$ & 0.09 & 0.1 & 0.23 & 0.36 & 0.49 \\ 
    & adj. $R^2$ & 0.06 & 0.04 & 0.14 & 0.26 & 0.36 \\ 
   \hline & Vital nodes & No & Yes & Yes & No & Yes \\ & Non-vital nodes & No & No & Yes & No & Yes \\ & Network activity & Yes & Yes & Yes & Yes & Yes \\ & Network statistics & No & No & No & Yes & Yes \\ \hline
\hline
\end{tabular}
\endgroup
\end{table}

%\input{table_hourly_VOL.tex}